\documentclass[conference]{IEEEtran}
\IEEEoverridecommandlockouts

% ── Packages ──────────────────────────────────────────────────────────────────
\usepackage{cite}
\usepackage{amsmath,amssymb,amsfonts}
\usepackage{algorithmic}
\usepackage{algorithm}
\usepackage{graphicx}
\usepackage{textcomp}
\usepackage{xcolor}
\usepackage{booktabs}
\usepackage{array}
\usepackage{url}
\usepackage{hyperref}
\usepackage{balance}
\usepackage{multirow}
\usepackage[T5]{fontenc}

\hypersetup{
    colorlinks=true,
    linkcolor=blue,
    citecolor=blue,
    urlcolor=blue
}

% ── Document ──────────────────────────────────────────────────────────────────
\begin{document}

\title{A Hybrid Smishing Detection Framework: Integrating Machine Learning,
Rule-Based Filtering, and Domain Reputation Analysis}



% \author{
%     \IEEEauthorblockN{
%     Duy Hoang Nguyen\IEEEauthorrefmark{1}\IEEEauthorrefmark{2}\IEEEauthorrefmark{3}\IEEEauthorrefmark{4},
%     Hai Nguyen-Nam Tran\IEEEauthorrefmark{1}\IEEEauthorrefmark{2}\IEEEauthorrefmark{3}\IEEEauthorrefmark{4},
%     Hai Thanh Tran\IEEEauthorrefmark{1}\IEEEauthorrefmark{2}\IEEEauthorrefmark{3}\IEEEauthorrefmark{4},\\
%     Quan Anh Vo\IEEEauthorrefmark{1}\IEEEauthorrefmark{2}\IEEEauthorrefmark{3}\IEEEauthorrefmark{4},
%     Vu Truong Cao\IEEEauthorrefmark{1}\IEEEauthorrefmark{2}\IEEEauthorrefmark{3}\IEEEauthorrefmark{4},
%     Tin Van Huynh\IEEEauthorrefmark{1}\IEEEauthorrefmark{2}\IEEEauthorrefmark{3}\IEEEauthorrefmark{5}
% }
%     \IEEEauthorblockA{
%         \IEEEauthorrefmark{1}Faculty of Information Science and Engineering\\
%         \IEEEauthorrefmark{2}University of Information Technology, Ho Chi Minh City, Vietnam\\
%         \IEEEauthorrefmark{3}Vietnam National University, Ho Chi Minh City, Vietnam\\
%         \IEEEauthorrefmark{4}\texttt{\{22520327, 22520391, 22520392, 22521192, 22521686\}@gm.uit.edu.vn}\\
%         \IEEEauthorrefmark{5}\texttt{tinhv@uit.edu.vn}
%     }
% }



\author{
    \IEEEauthorblockN{Anonymous Authors}
}

\maketitle

% ── Abstract ──────────────────────────────────────────────────────────────────
\begin{abstract}
The rise of SMS phishing (smishing) attacks poses a significant threat to
mobile security, particularly for Vietnamese users whose linguistic environment
is inadequately served by existing international solutions.
This paper makes three primary contributions.
First, we construct a Vietnamese smishing dataset comprising 2,600
labeled SMS messages, collected from real-world devices and
community-sourced reports.
Second, we propose a Hybrid Smishing Detection Framework featuring a
four-layer linguistic preprocessing pipeline designed to handle
Vietnamese text obfuscation, including teencode, Leetspeak, and broken Telex
input.
Third, we integrate this pipeline with a tree-based ML classifier and a
Hybrid Decision Logic that incorporates Domain Verification.
On a held-out test set of 520 samples, the XGBoost-based system achieves an
overall accuracy of 97.5\% and a smishing recall of 78\%, outperforming
single-model baselines in the precision-recall trade-off relevant to
cyber-security deployment.
\end{abstract}

\begin{IEEEkeywords}
Smishing, SMS Phishing, Machine Learning, Vietnamese NLP,
Linguistic Feature Extraction, Hybrid Architecture
\end{IEEEkeywords}

% ── 1. INTRODUCTION ───────────────────────────────────────────────────────────
\section{Introduction}

In the modern digital era, SMS remains one of the most reliable communication
channels, with an open rate and near-immediate readership that far surpass those of other channels such as email.
This trust has made it a prime vector for \textit{smishing} (SMS Phishing) - a
form of social engineering where attackers deceive victims into revealing
sensitive information or installing malware. In Vietnam, attackers routinely
impersonate banks, e-commerce platforms, and government agencies, often
exploiting Brandname abuse to bypass traditional filters, causing significant
financial and data losses for unsuspecting users.

Existing defenses remain limited: static blacklists and simple keyword
matching fail against obfuscated text (teencode, Leetspeak, mixed-character
URLs) and dynamic URL shorteners. Advanced ML pipelines from the international
literature are furthermore optimized for English; applying them to Vietnamese
yields suboptimal results due to diacritical marks, tone-dependent semantics,
and local obfuscation conventions. To the best of our knowledge, no
comprehensive system currently adapts these architectures specifically for
the Vietnamese linguistic and telecommunication landscape.

We define the smishing detection task as a binary text classification problem.
Given an incoming SMS message $T$ with sender type $S$, the system must assign
a label $L \in \{\textit{Smishing},\, \textit{Legitimate}\}$. Formally:

\begin{itemize}
    \item \textbf{Input:} A tuple $(T, S)$, where $T$ is the raw SMS message
          body and $S$ denotes the sender type (e.g. Brandname, Phone Number,
          Unknown).
    \item \textbf{Output:} A label $L$, a confidence score
          $\text{Score} \in [0,1]$, and a natural-language explanation of the
          decision.
\end{itemize}

\begin{figure}[htbp]
    \centering
    \includegraphics[width=\columnwidth]{fig1_examples.jpg}
    \caption{Exambples for each type of SMS.}
    \label{fig:domain}
\end{figure}

To address the challenges above, we propose a Hybrid Real-Time Smishing
Detection System tailored for Vietnamese. Our key contributions are:
\begin{enumerate}
    \item \textbf{Dataset:} A composite Vietnamese dataset of 2,600 labeled
          SMS samples built via dual-channel acquisition and a five-annotator
          majority-vote protocol.
    \item \textbf{Framework:} A four-layer linguistic preprocessing pipeline
          for Vietnamese, combined with an ML classifier and deterministic
          Domain Verification integrated through a Hybrid Decision Logic.
    \item \textbf{Experiments:} Comparative evaluation of Decision Tree,
          Random Forest, and XGBoost demonstrating our hybrid approach
          outperforms TF-IDF baselines in the precision-recall trade-off
          critical for cyber-security deployment.
\end{enumerate}
% ── 2. RELATED WORKS ──────────────────────────────────────────────────────────
\section{Related Works}
\label{sec:related}

\subsection{Related Methods}

The landscape of smishing detection has evolved from simple rule-based filters
toward increasingly sophisticated machine learning and deep learning approaches.
Broad surveys confirm that machine learning is among the most promising
anti-smishing methods, yet continuous evaluation of algorithms is required as
attacker tactics evolve~\cite{mahmood2023review}. Within this space, deep
learning architectures have attracted significant attention: Mahmud
\textit{et al.}~\cite{mahmud2024hybrid} propose a hybrid model combining
Bidirectional Gated Recurrent Units (Bi-GRUs) with Convolutional Neural
Networks (CNNs) for smishing classification, while Mehmood
\textit{et al.}~\cite{mehmood2024smishing} aggregate CNNs and Long Short-Term
Memory (LSTM) networks to automatically extract features from SMS text and
reduce false-positive rates. Beyond text-only input, Shinde
\textit{et al.}~\cite{shinde2024ocr} expand the feature space using optical
character recognition (OCR) to detect scams embedded in image-based messages,
combining unsupervised and deep semi-supervised learning.

Despite their expressive power, deep architectures demand large labeled
corpora and sacrifice interpretability-two critical constraints in our
setting. Mishra and Soni~\cite{mishra2023dsmishsms} demonstrate that a
Defense-in-Depth hybrid coupling a Domain Checking Phase (URL validation via
search-engine signatures) with an SMS Classification Phase based on five
hand-crafted features achieves 97.93\% accuracy while remaining fully
interpretable. Goel \textit{et al.}~\cite{goel2024smishing} further show
that explicit text normalization before classification yields 96.2\% accuracy
on an English dataset; we extend this normalization philosophy to Vietnamese,
where tone-mark interactions, Telex artifacts, and Leet obfuscation demand a
multi-layer pipeline rather than a single normalization pass. Mishra and Soni
also found that 74.22\% of smishing messages contain character substitutions
(e.g. ``P@ytm''$\to$``Paytm'') to bypass keyword filters, a finding
corroborated by Li and Zeng~\cite{li2021leet} who document how such
substitutions systematically evade dictionary-based detectors-a challenge
amplified in Vietnamese due to diacritical marks. Crucially, Tuan
\textit{et al.}~\cite{tuan2023vietnamese} evaluate SVM, Naïve Bayes, Random
Forests, CNN, and LSTM specifically on Vietnamese SMS, noting measurable
performance differences when processing text with and without diacritics and
confirming that language-specific treatment is indispensable. Building on
these findings, we adopt XGBoost~\cite{chen2016xgboost}-a scalable
tree-boosting system with a sparsity-aware split-finding algorithm and
second-order gradient statistics-as our core classifier, since the
28-dimensional dense feature vector produced by our linguistic pipeline
constitutes exactly the kind of structured, heterogeneous input for which
gradient-boosted trees excel over both shallow models and sparse TF-IDF
representations, while remaining interpretable and efficient for real-time
deployment.

\subsection{Related Datasets}

A persistent challenge in smishing research is the scarcity of labeled,
publicly available datasets that reflect real-world attack patterns. Almeida
\textit{et al.}~\cite{almeida2011sms} introduced the UCI SMS Spam Collection,
comprising 5,574 English-language messages, which became the de facto
benchmark for SMS spam research; however, it does not capture the linguistic
properties of other languages. More recently, Timko and
Rahman~\cite{timko2024smishtank} compiled a large community-sourced smishing
corpus from Smishtank.com, providing detailed metadata including sender
information, message bodies, and referenced brands-advancing the field with
richer annotation but remaining English-centric.

On the Vietnamese side, Nguyen \textit{et al.}~\cite{nhi2024vilexnorm}
developed ViLexNorm, a lexical normalization corpus of annotated Vietnamese
social media sentence pairs designed to address intended and unintended
linguistic variances such as teencode and Leetspeak. While ViLexNorm provides
a valuable normalization resource, no publicly available Vietnamese-specific
smishing dataset exists. Our work addresses this gap by constructing a corpus
of 2,600 labeled SMS messages through a dual-channel acquisition strategy,
enabling evaluation under Vietnamese obfuscation conventions and
Brandname-abuse patterns that no existing public dataset covers.

% ── 3. DATA COLLECTION ────────────────────────────────────────────────────────
\section{Data Collection}
\label{sec:data}

Due to the scarcity of standardized Vietnamese smishing datasets, we
constructed a composite corpus through a dual-channel acquisition strategy.

\paragraph{Manual Transcription from Public Sources.}

The first channel harvested confirmed smishing messages frequently reported
as screenshots on Vietnamese cyber-security forums and anti-scam community
groups. We performed manual transcription rather than OCR to accurately
preserve linguistic nuances-teencode, special character obfuscations, and
malicious URLs - that automated recognition tools often corrupt. This yielded
a high-quality set of distinct smishing signatures targeting banking and
e-commerce sectors.

\subsection{Real-World Data Acquisition via Mobile Utility}

The second channel collected real-world SMS data-legitimate conversations,
OTPs, and carrier advertisements-from real Android devices using
\textit{SMS Backup \& Restore}, with explicit consent from participants.
This application exports local message repositories into a structured XML
format, allowing efficient extraction of message body, timestamp, and
sender information.

\paragraph{Real-World Data Acquisition via Mobile Utility.}

Following aggregation, cleaning, and normalization, the final dataset
comprises 2,600 unique rows (Table~\ref{tab:dataset}). Labels were assigned
by five annotators using pre-defined guidelines: a message is
\textit{Smishing} if it contains (1) a URL leading to a non-whitelisted
or suspicious domain, (2) impersonation of a trusted entity, (3) financial
pressure or urgency combined with a request for personal information, or
(4) unrealistic recruitment offers. Final labels were decided by majority
vote; unresolved samples were discarded to maintain dataset purity.

\begin{table}[htbp]
    \caption{Dataset Class Distribution}
    \label{tab:dataset}
    \centering
    \begin{tabular}{lcc}
        \toprule
        \textbf{Class} & \textbf{Count} & \textbf{Percentage} \\
        \midrule
        Legitimate (Ham) & 2,353 & 90.5\% \\
        Smishing (Spam)  & 247   &  9.5\% \\
        \midrule
        \textbf{Total}   & \textbf{2,600} & \textbf{100\%}  \\
        \bottomrule
    \end{tabular}
\end{table}

% ── 4. DATA PREPROCESSING ─────────────────────────────────────────────────────
\section{Data Preprocessing}
\label{sec:preprocess}

SMS data contains significant noise and lacks structural uniformity. In spam
messages, obfuscation techniques-inserting random spaces within URL
segments, applying Leetspeak, or mixing characters-actively bypass
keyword-based filters. A single-pass normalization approach fails to address
these interacting noise sources, motivating our four-layer design. Each layer
targets a distinct noise category; their outputs also serve as discriminative
classifier features.

\paragraph{Layer 1 - Entity Masking.}
Specific entities are replaced with standardized tokens (e.g. \texttt{<URL>},
\texttt{<PHONE>}, \texttt{<MONEY>}) using an OrderedDict to enforce
conflict-free processing order. The pipeline executes in the following order:
\begin{itemize}
    \item \textbf{Platform Links:} Zalo and Telegram links (\texttt{<LINK>}).
    \item \textbf{Emails, URLs:} Standard email regex; custom logic for
          obfuscated malicious links (\texttt{<EMAIL>}, \texttt{<URL>}).
    \item \textbf{Bank Accounts, Telecommunications:} Context-aware
          extraction for hotlines, landlines, and mobile numbers.
    \item \textbf{Datetime, Monetary Values, Codes \& OTPs:}
          \texttt{<TIME>}, \texttt{<MONEY>}, \texttt{<CODE>}.
    \item \textbf{Anti-Evasion URL Detection:} Detects intentional spacing
          in domains (e.g. bit.\,ly/abc), and protocol-agnostic
          domains matching high-risk TLD patterns (.xyz, .top, .moe).
\end{itemize}

\paragraph{Layer 2 - Text Normalization and Leetspeak Decoding.}
The \textit{TextNormalizer} standardizes the character set and decodes
obfuscations via a weighted scoring system validated against a Vietnamese
dictionary. Three levels of obfuscation are recognized:
\begin{itemize}
    \item \textbf{Symbol Leet} (weight 1.2--1.5): e.g. @$\to$a,
          !$\to$i. Highly indicative of spam.
    \item \textbf{Visual Leet} (weight 1.0): e.g. 0$\to$o,
          4$\to$a (\textit{h0tro} $\to$ \textit{ho tro}).
    \item \textbf{Teencode} (weight 0.3): e.g. \textit{ko}$\to$\textit{khong}.
\end{itemize}
A placeholder mechanism protects Layer~1 tokens from corruption during
separator removal.

\paragraph{Layer 3 - Whitelist Filtering.}
Valid out-of-dictionary tokens - brand names (Shopee, VCB), technical jargon
(OTP, 4G), and SMS abbreviations (\textit{lh}, \textit{sdt}) - are exempted
from further processing to prevent erroneous spell-check corrections
(e.g. ``VCB'' $\to$ ``viec'').

\paragraph{Layer 4 - Spell Check and Anomaly Detection.}
Remaining Out-Of-Vocabulary (OOV) tokens are analyzed by a Dual-Lookup
Mechanism (full-diacritics dictionary + accent-free shadow dictionary).
Tokens failing both checks are tagged as one of four error categories used
as classifier features: \textit{Telex Artifacts} (aa, aw,
dd), \textit{Run-On Words} (bisection heuristic, e.g.
\textit{vui\,+\,long}), \textit{Gibberish Strings} (violating Vietnamese
phonotactics), or \textit{Character Repetition} (3+ identical characters,
e.g. \textit{HOTTT}). A Data Capacity Whitelist proactively exempts
telco tokens (e.g. \texttt{3gb}, \texttt{12mbps}) that are technically OOV
but valid in SMS contexts.

% ── 5. PROPOSED SYSTEM ────────────────────────────────────────────────────────
\section{Proposed System Methodology}
\label{sec:method}

The framework combines a Machine Learning scorer, a Domain Verification
module, and a Hybrid Decision Logic. A purely statistical model is vulnerable
to adversarial evasion; a purely rule-based system cannot generalize to novel
attack patterns. The hybrid architecture addresses both shortcomings.

\subsection{Machine Learning Model}
\label{subsec:ml}

The raw SMS $T$ is passed through the 4-Layer Pipeline
(Section~\ref{sec:preprocess}) to generate a 27-dimensional linguistic vector
$F_\text{vec}$. Unlike sparse TF-IDF matrices, this dense vector directly
encodes obfuscation indicators prevalent in Vietnamese fraud. The Sender
Type $S$ (Brandname, Phone Number, etc.) is label-encoded and appended:
\begin{equation}
    V = \text{COMBINE}(F_\text{vec},\, S) \in \mathbb{R}^{28}
\end{equation}
The pre-trained XGBoost model $M_{AI}$ produces a probability score:
\begin{equation}
    P_{AI} = M_{AI}.\text{predict\_proba}(V), \quad P_{AI} \in [0, 1]
\end{equation}
A score near~1 reflects high obfuscation and urgency cues; near~0 indicates
legitimate patterns. $P_{AI}$ feeds the Hybrid Decision Logic.

\subsection{Domain Verification}
\label{subsec:domain}

The Domain Verification module provides deterministic, AI-independent
verdicts for URL-bearing messages (Fig.~\ref{fig:domain}). The brand context
$B$ and all URLs are extracted; \textbf{tldextract} isolates the registered
domain $D$ to prevent subdomain spoofing (e.g. distinguishing
\textit{vietcombank.ngrok.io} as belonging to \textit{ngrok.io}).
If $D \in W_\text{static}$ (a pre-compiled whitelist of major Vietnamese
banks, telcos, and global platforms), the domain is flagged Legitimate
($\text{Risk}{=}{-1.0}$). Otherwise, two-level dynamic verification triggers:
Level~1 checks whether $D$ appears in the top-5 search results for the brand
name; Level~2 scans the site's source for internal links pointing to verified
domains. An unreachable, redirect-heavy, or unverifiable site yields
$\text{Risk}{=}1.0$.

\begin{figure}[htbp]
    \centering
    \includegraphics[width=\columnwidth]{fig2_domain.jpg}
    \caption{Architecture of the Domain Verification Phase.}
    \label{fig:domain}
\end{figure}

\subsection{Hybrid Decision Logic}
\label{subsec:hybrid}

The Hybrid Decision Logic applies a three-tier waterfall
(Algorithm~\ref{alg:hybrid}; system architecture in Fig.~\ref{fig:arch}).

\noindent\textbf{Tier~1 - Domain-Driven Verdict:} If $\text{Risk} \geq 0.8$,
the message is immediately classified as Smishing. For trusted domains
($\text{Risk}{=}{-1.0}$), user-generated-content platforms undergo a secondary
AI check ($P_{AI}{>}0.65$).

\noindent\textbf{Tier~2 - Context-Aware AI Verdict:} When no URL is present
or domain status is inconclusive, $P_{AI}$ is evaluated alongside context
flags. Conversational messages override high AI risk; messages with danger
keywords from unknown senders force a Smishing verdict.

\noindent\textbf{Tier~3 - Explainable Output:} Every decision yields a
tuple $(\text{Label},\,\text{Score},\,\text{Decision Phase},\,\text{Reasoning})$
for user and analyst transparency.

\begin{figure}[t]
    \centering
    \includegraphics[width=\columnwidth]{fig3_architecture.jpg}
    \caption{Architecture of the proposed Hybrid Smishing Detection System.}
    \label{fig:arch}
\end{figure}

\begin{algorithm}[t]
\caption{Hybrid Smishing Detection Algorithm}
\label{alg:hybrid}
\begin{algorithmic}[1]
\REQUIRE Incoming SMS $T$, Sender Type $S$, Thresholds $T_{AI}$, $T_\text{high}$
\ENSURE Label (\textit{Smishing} / \textit{Legitimate})
\STATE $F_\text{vec} \leftarrow \text{ExtractFeatures}(T)$ \COMMENT{Layers 1--4}
\STATE $V \leftarrow \text{Combine}(F_\text{vec}, S)$
\STATE $P_{AI} \leftarrow M_{AI}.\text{predict\_proba}(V)$
\STATE $\text{Risk} \leftarrow 0$
\IF{$T$ contains URL}
    \STATE $\text{Risk} \leftarrow \text{VerifyDomain}(T)$
\ENDIF
\IF{$\text{Risk} \geq T_\text{high}$}
    \RETURN \textit{Smishing} \COMMENT{Tier 1: High-Risk Domain}
\ELSIF{$\text{Risk} = -1.0$}
    \IF{$P_{AI} > 0.65$}
        \RETURN \textit{Smishing (Warning)} \COMMENT{UGC Exception}
    \ELSE
        \RETURN \textit{Legitimate}
    \ENDIF
\ELSE
    \IF{$P_{AI} \geq T_{AI}$}
        \RETURN \textit{Smishing}
    \ELSE
        \RETURN \textit{Legitimate}
    \ENDIF
\ENDIF
\end{algorithmic}
\end{algorithm}

% ── 6. EXPERIMENTS ────────────────────────────────────────────────────────────
\section{Experiments}
\label{sec:eval}
 
\subsection{Methods and Experimental Setup}
 
The 2,600-message dataset (90\% Legitimate, 10\% Smishing) was split 80/20
using stratified sampling (\texttt{random\_state=42, stratify=y}), yielding a test set
of 520 messages (465 Legitimate, 55 Smishing). We adopt a two-level evaluation
strategy: (i)~three standalone classifiers evaluated on the full 520-sample
test set using 28-dimensional features, and (ii)~the complete Hybrid System
evaluated on 100 stratified messages with per-error categorization.
 
\subsection{Evaluation Metrics}
 
All Precision, Recall, and F1 values refer to the \textit{Smishing} minority
class, the operationally relevant class in cyber-security deployment. Metrics
are reported on $[0,1]$. Rather than a fixed threshold of 0.5, each
classifier's optimal threshold $\tau^{*}$ is selected as the maximum-F1
operating point on the Precision-Recall curve:
\[
  \tau^{*} = \arg\max_{\tau}
    \frac{2 \cdot P(\tau)\cdot R(\tau)}{P(\tau)+R(\tau)}.
\]
The Smishing F1-Score is the primary comparison metric, as overall Accuracy
can be misleading under a 90/10 class imbalance~\cite{he2009imbalanced}.
 
\subsection{Experimental Results}
\label{subsec:comparison}
 
We compare the 4-Layer Pipeline against a TF-IDF Baseline (5,000
unigrams/bigrams augmented with \textit{has\_url}, \textit{has\_phone},
\textit{sender\_type}; fixed threshold 0.5). The results are presented in Table~\ref{tab:compare_baseline}.
 
\begin{table}[htbp]
  \caption{TF-IDF Baseline vs.\ Proposed 4-Layer Linguistic Pipeline}
  \label{tab:compare_baseline}
  \centering
  \resizebox{\columnwidth}{!}{%
  \begin{tabular}{llccccc}
    \toprule
    \textbf{Feature Set} & \textbf{Model}
      & \textbf{Acc.} & \textbf{Prec.} & \textbf{Recall}
      & \textbf{F1} & \textbf{ROC-AUC} \\
    \midrule
    \multirow{3}{*}{TF-IDF Baseline ($\tau{=}0.5$)}
      & Decision Tree  & 0.9231 & 0.6000 & 0.8182 & 0.6923 & 0.8831 \\
      & Random Forest  & 0.9673 & 0.9750 & 0.7091 & 0.8211 & 0.9844 \\
      & XGBoost        & 0.9673 & 0.8800 & 0.8000 & 0.8381 & 0.9696 \\
    \midrule
    \multirow{3}{*}{4-Layer Linguistic ($\tau{=}\tau^{*}$)}
      & Decision Tree  & 0.9519 & 0.7679 & 0.7818 & 0.7748 & 0.9050 \\
      & Random Forest  & 0.9654 & 0.9744 & 0.6909 & 0.8085 & 0.9730 \\
      & \textbf{XGBoost}
        & \textbf{0.9750} & \textbf{0.9773} & \textbf{0.7818}
        & \textbf{0.8687} & \textbf{0.9659} \\
    \bottomrule
  \end{tabular}%
  }
  \\[4pt]
  {\small Prec., Recall, F1 are for the Smishing class.}
\end{table}
 
The linguistic pipeline yields the most pronounced gains for XGBoost:
Smishing F1 rises from 0.838 to \textbf{0.869} ($+$3.1~pp), and-crucially
for operational deployment-Precision rises from 0.88 to \textbf{0.977}
($+$9.7~pp), meaning fewer than 2.3\% of flagged messages are false alarms.
Decision Tree benefits most in F1 ($+$8.3~pp), confirming that hand-crafted
obfuscation indicators carry discriminative signal that sparse TF-IDF cannot
capture for leet-encoded Vietnamese text.
 
\subsection{Standalone Model Results}
\label{subsec:analysis}
 
Table~\ref{tab:ml_results} reports full results at optimal thresholds.
 
\begin{table}[htbp]
  \caption{Smishing-Class Metrics at Optimal Threshold
           (4-Layer Linguistic Features; test set $n{=}520$)}
  \label{tab:ml_results}
  \centering
  \resizebox{\columnwidth}{!}{%
  \begin{tabular}{lccccccc}
    \toprule
    \textbf{Model} & $\tau^{*}$
      & \textbf{Acc.} & \textbf{Prec.} & \textbf{Recall} & \textbf{F1}
      & \textbf{TP} & \textbf{FP} \\
    \midrule
    Decision Tree        & 0.8942 & 0.9519 & 0.7679 & 0.7818 & 0.7748 & 43 & 13 \\
    Random Forest        & 0.6842 & 0.9654 & 0.9744 & 0.6909 & 0.8085 & 38 &  1 \\
    \textbf{XGBoost}     & \textbf{0.8336}
      & \textbf{0.9750} & \textbf{0.9773} & \textbf{0.7818}
      & \textbf{0.8687} & \textbf{43} & \textbf{1} \\
    \bottomrule
  \end{tabular}
  }
  \\[4pt]
  {\small FN $= 55 - \text{TP}$; TN $= 465 - \text{FP}$.}
\end{table}

Across all three models, there is a trade-off between Precision and Recall
that reflects a common challenge in imbalanced classification. Decision Tree
achieves the highest Recall (0.782) but incurs 13 false positives, while
Random Forest holds Precision at 0.974 at the cost of missing 17 smishing
messages (Recall $= 0.691$). XGBoost balances both criteria, matching
Decision Tree's Recall while reducing false positives to 1 - yielding the
highest Smishing F1 (0.869) and overall Accuracy (0.975) among the
three classifiers.
 
% ── 7. RESULT ANALYSIS ────────────────────────────────────────────────────────
\section{Result Analysis}
\label{sec:result_analysis}
 
\subsection{Standalone Model Analysis}
 
\subsubsection*{Decision Tree}
The Decision Tree detects 43 of 55 smishing messages
(Recall~$= 0.78$), but at the cost of 13 false positives - roughly
1 in 4 flagged messages is actually benign (Precision~$= 0.768$), a
high false-alarm rate for production use. Overall accuracy is 95.2\%.
 
\subsubsection*{Random Forest}
Random Forest produces only 1 false positive across 465 legitimate
messages (Precision~$= 0.974$), but misses 17 smishing messages
(Recall~$= 0.69$; 31\% miss rate). In a security context, missing nearly
one in three attacks is a notable drawback. Random Forest may be better
suited to scenarios where false alarms carry higher cost than missed attacks.
 
\subsubsection*{XGBoost}
XGBoost achieves a balanced trade-off for cyber-security deployment. It
detects the same 43 of 55 smishing messages as Decision Tree
while producing only 1 false positive - catching 5 more
attacks than Random Forest ($+$8.3~pp Recall) at equal false-alarm cost.
Its Smishing F1 of 0.8687 is the highest among the three, and its
overall Accuracy of 0.9750 reflects consistent performance on both
classes. XGBoost is therefore selected as the core classifier of the
Hybrid Detection Framework.
 
\subsection{Error Analysis of the Full Hybrid System}
\label{subsec:error}
 
The complete Hybrid System ($\tau \approx 0.46$) was applied to
100 stratified messages. Of these, 87 were correctly
classified, yielding 3 False Negatives (FN) and
10 False Positives (FP), summarized in Table~\ref{tab:error}.
 
\begin{table}[htbp]
  \caption{Error Analysis on 100 Sampled Messages (Hybrid System)}
  \label{tab:error}
  \centering
  \begin{tabular}{llc}
    \toprule
    \textbf{Type} & \textbf{Category} & \textbf{Count} \\
    \midrule
    FN & No-URL / low-obfuscation smishing below AI threshold & 3 \\
    \midrule
    FP & AI model over-triggered on transactional messages   & 5 \\
    FP & Legitimate domain absent from Domain Risk whitelist & 5 \\
    \midrule
    \multicolumn{2}{l}{\textbf{Total errors}} & \textbf{13} \\
    \bottomrule
  \end{tabular}
\end{table}
 
\subsubsection*{False Negatives}
All three FNs are smishing messages where neither the AI model nor the Domain
Risk module exceeded their respective thresholds. Three distinct failure
patterns were identified:
\begin{enumerate}
  \item \textbf{Authority impersonation without a URL.} A message
        impersonating \textit{C\'{a}nh s\'{a}t Giao th\^{o}ng}-containing
        only a hotline number-received $P_{AI}{=}0.20$ because it is
        written in clean, formal Vietnamese that lacks the obfuscation cues
        the model relies on. The Domain Risk module had no URL to evaluate.
  \item \textbf{Telegram social-engineering link with minimal obfuscation.}
        A message linking to \texttt{t.me/hn\_silent} was scored
        $\text{Risk}{=}0.0$ because \texttt{t.me} is whitelisted as a known
        messaging platform, and the AI gave $P_{AI}{=}0.09$ as the text
        resembles a social invitation rather than a financial scam.
  \item \textbf{Leet-speak phishing URL just below the decision boundary.}
        A VNeID impersonation with leet-encoded text obtained
        $P_{AI}{=}0.45$ - one point below $\tau^{*}{\approx}0.46$-while
        the Domain Risk module returned $r{=}0.0$ for
        \texttt{vne1d-update.t0p}, finding no conclusive blacklist evidence.
\end{enumerate}

These cases motivate: (a)~a recall-biased rule for URL-free
government/law-enforcement impersonation, and (b)~a pattern-based leet-URL
detector for domains whose structure matches known brand names.
 
\subsubsection*{False Positives}
Both FP clusters arise from the system's inability to distinguish
surface-level similarity from genuine fraud intent. Two distinct failure
patterns were identified:
\begin{enumerate}
  \item \textbf{AI over-triggered on transactional messages.} Five legitimate
        confirmation messages (restaurant reservations, taxi bookings,
        e-commerce notifications) were flagged because they share structural
        features with smishing: brevity, a specific time reference, and an
        entity name (e.g. a reservation confirmation scored
        $P_{AI}{=}0.94$). The classifier has learned co-occurrence of
        \{short text, temporal marker, named entity\} as a fraud indicator
        without encoding the semantic distinction between transactional
        confirmations and fraudulent solicitations. Remedy: augment training
        data with hard-negative transactional messages and introduce a
        dedicated \textit{transactional intent} feature.
  \item \textbf{Legitimate domain absent from the whitelist.} Five messages
        were escalated by the Domain Risk module because their domains were
        absent from the static whitelist and the live-search verifier returned
        inconclusive results ($\text{Risk}{=}1.0$). Affected domains include
        \texttt{thuedientu.gdt.gov.vn} (official tax portal),
        \texttt{ghn.vn} (GHN Express), and \texttt{shopeepay.vn}.
        Remedy: extend the whitelist to cover all \texttt{*.gov.vn} and
        \texttt{*.edu.vn} subdomains, and apply a low-risk prior for
        government and academic TLDs in the live-search verifier.
\end{enumerate}
These cases motivate: (a)~richer semantic features to separate transactional
from fraudulent intent in the AI tier, and (b)~a broader, programmatically
maintained domain whitelist that covers official government and
logistics namespaces.
 
% ── 8. CONCLUSION ─────────────────────────────────────────────────────────────
\section{Conclusion and Future Work}
\label{sec:conclusion}

This paper presented a Hybrid Smishing Detection Framework purpose-built for
the Vietnamese linguistic environment and made three primary contributions.
\textbf{First}, we constructed a Vietnamese smishing dataset of 2,600 labeled
SMS messages via a dual-channel acquisition strategy - manual transcription
from community anti-scam sources and real-device extraction via mobile backup
utilities - with a five-annotator majority-vote labeling protocol, providing a
reproducible benchmark with a realistic 90/10 class imbalance.
\textbf{Second}, we proposed a four-layer linguistic preprocessing pipeline
(\textit{Entity Masking} $\to$ \textit{Leet Normalization} $\to$
\textit{Whitelist Filtering} $\to$ \textit{Anomaly Detection}) converting raw,
obfuscated Vietnamese SMS into a dense 28-dimensional feature vector. Compared
to a TF-IDF baseline, the pipeline raised Smishing F1 from 0.838 to
0.869 ($+$3.1~pp) and Precision from 0.88 to 0.977
($+$9.7~pp), reducing false alarms to below 2.3\%.
\textbf{Third}, we integrated the classifier into a three-tier Hybrid Decision
Logic (Domain-Driven Verdict, Context-Aware AI Verdict, Explainable Output)
that allows high-confidence domain signals to override AI scores in both
directions. On a 100-message real-world evaluation the system correctly
classified 87\% of samples. In aggregate, the framework achieves
97.5\% overall accuracy, 78.2\% Smishing Recall, and
97.7\% Smishing Precision on the held-out test set, demonstrating that
language-specific obfuscation features provide strong discriminative signal for
tree-based classifiers while maintaining real-time mobile inference speed.

Four limitations constrain the current system. \textbf{Dataset scale:} 2,600
samples with only 247 positives limits coverage of rare attack sub-types such
as crypto fraud and OTP relay attacks. \textbf{Static rules:} hand-crafted
pipeline features require manual maintenance as attackers adapt, with whitelist
blind spots confirmed by the error analysis. \textbf{URL dependence:} all three
false negatives were URL-free messages where the 28-dimensional vector cannot
capture semantic intent for clean-text authority-impersonation. \textbf{SMS
scope:} the pipeline does not model Zalo or Telegram channels where image-based
content and group-chat dynamics introduce unseen evasion patterns.

Three future directions address these gaps. \textbf{LLM-Driven Dataset
Expansion} will grow the corpus to 10,000$+$ samples using GPT-4, Claude, and
Gemini as controlled generators of hard negatives and adversarial smishing
examples, each annotated with a \textit{Rationale} explaining the
classification. \textbf{Neural Text Normalization} will replace brittle regex
rules with a BARTpho-based~\cite{nguyen2022bartpho} Seq2Seq normalizer trained
on ViLexNorm~\cite{nhi2024vilexnorm} to map all obfuscation variants of a token
to a single canonical form. \textbf{Rationale-Aware Knowledge Distillation}
will train a compact ViSoBERT~\cite{nguyen2023visobert} Student via
$\mathcal{L}_\text{total} = \alpha\,\mathcal{L}_\text{Hard}
+ \beta\,\mathcal{L}_\text{Soft} + \gamma\,\mathcal{L}_\text{Rationale}$,
where $\mathcal{L}_\text{Hard}$ provides ground-truth supervision,
$\mathcal{L}_\text{Soft}$ transfers Teacher uncertainty on borderline cases, and
$\mathcal{L}_\text{Rationale}$ (Cosine Similarity) aligns the Student's hidden
state with the Teacher's semantic reasoning - targeting $>$ 90\% Smishing F1 and
robustness against both identified failure modes.

% \section*{Acknowledgment}
% This research was supported by The VNUHCM-University of Information Technology's Scientific Research Support Fund.

% ── REFERENCES ────────────────────────────────────────────────────────────────
\balance
\begin{thebibliography}{00}

\bibitem{mishra2023dsmishsms}
S. Mishra and D. Soni,
``DSmishSMS-A System to Detect Smishing SMS,''
\textit{Neural Computing and Applications},
vol.~35, pp.~4975--4992, 2023.

\bibitem{li2021leet}
W. Li and J. Zeng,
``Leet Usage and Its Effect on Password Security,''
\textit{IEEE Transactions on Information Forensics and Security},
vol.~16, pp.~2130--2143, 2021.

\bibitem{chen2016xgboost}
T. Chen and C. Guestrin,
``XGBoost: A Scalable Tree Boosting System,''
in \textit{Proc. 22nd ACM SIGKDD Int. Conf. Knowledge Discovery and
Data Mining (KDD)}, San Francisco, CA, Aug.\ 2016, pp.\ 785--794.
DOI: 10.1145/2939672.2939785.

\bibitem{almeida2011sms}
T.~A. Almeida, J.~M.~G. Hidalgo, and A. Yamakami,
``Contributions to the Study of SMS Spam Filtering: New Collection and
Results,''
in \textit{Proc. 11th ACM Symp. Document Engineering (DOCENG)},
Mountain View, CA, Sep.\ 2011, pp.\ 259--263.
DOI: 10.1145/2034691.2034742.

\bibitem{goel2024smishing}
D.~Goel, H. Ahmad, A.~K. Jain, and N.~K. Goel,
``Machine Learning Driven Smishing Detection Framework for Mobile Security,''
\textit{arXiv preprint arXiv:2412.09641}, Dec.\ 2024.

\bibitem{mahmood2023review}
A.~R. Mahmood and S.~M. Hameed,
``Review of Smishing Detection via Machine Learning,''
\textit{Iraqi Journal of Science},
pp.~4244--4259, 2023.
DOI: 10.24996/ijs.2023.64.8.42.

\bibitem{mahmud2024hybrid}
T. Mahmud, M.~A.~H. Prince, M.~H. Ali, M.~S. Hossain, and K. Andersson,
``Enhancing Cybersecurity: Hybrid Deep Learning Approaches to Smishing
Attack Detection,''
\textit{Systems}, vol.~12, p.~490, 2024.
DOI: 10.3390/systems12110490.

\bibitem{mehmood2024smishing}
M.~K. Mehmood, H. Arshad, M. Alawida, and A. Mehmood,
``Enhancing Smishing Detection: A Deep Learning Approach for Improved
Accuracy and Reduced False Positives,''
\textit{IEEE Access}, vol.~12, pp.~137176--137193, 2024.
DOI: 10.1109/access.2024.3463871.

\bibitem{shinde2024ocr}
A. Shinde, E.~Q. Shahra, S. Basurra, F. Saeed, A.~A. AlSewari,
and W.~A. Jabbar,
``SMS Scam Detection Application Based on Optical Character Recognition
for Image Data Using Unsupervised and Deep Semi-Supervised Learning,''
\textit{Sensors}, vol.~24, p.~6084, 2024.
DOI: 10.3390/s24186084.

\bibitem{timko2024smishtank}
D. Timko and M.~L. Rahman,
``Smishing Dataset I: Phishing SMS Dataset from Smishtank.com,''
in \textit{Proc. ACM CODASPY}, 2024.
DOI: 10.1145/3626232.3653282.

\bibitem{tuan2023vietnamese}
V.~M. Tuan, D.~T. Duong, and T.~Q. Anh,
``Proposing Appropriate SMS Spam Detection Approaches for Variations of
the Vietnamese Language,''
in \textit{Proc. 2023 RIVF Int. Conf. Computing and Communication
Technologies (RIVF)}, pp.~1--5, 2023.
DOI: 10.1109/rivf60135.2023.10471859.

\bibitem{he2009imbalanced}
H. He and E.~A. Garcia,
``Learning from Imbalanced Data,''
\textit{IEEE Transactions on Knowledge and Data Engineering},
vol.~21, no.~9, pp.~1263--1284, Sep.\ 2009.
DOI: 10.1109/TKDE.2008.239.

\bibitem{nguyen2022bartpho}
N.~L. Nguyen, V.-L. Tran, D.~Q. Nguyen, T.~H. Nguyen, T.~T. Nguyen,
T.~M. Le, and T.~H. Nguyen,
``BARTpho: Pre-trained Sequence-to-Sequence Models for Vietnamese,''
in \textit{Proc. Language Resources and Evaluation Conf. (LREC)},
Marseille, France, Jun.\ 2022, pp.~4673--4682.

\bibitem{nhi2024vilexnorm}
T.-N. Nguyen, T.-P. Le, and K.~V. Nguyen,
``ViLexNorm: A Lexical Normalization Corpus for Vietnamese Social Media
Text,''
\textit{arXiv preprint arXiv:2401.16403}, 2024.
DOI: 10.48550/arxiv.2401.16403.

\bibitem{nguyen2023visobert}
N.~L. Nguyen, T.~N. Nguyen, K.~V. Nguyen, and N.~L.-T. Nguyen,
``ViSoBERT: A Pre-Trained Language Model for Vietnamese Social Media
Text Processing,''
in \textit{Proc. 2023 Conf. Empirical Methods in Natural Language
Processing (EMNLP)}, Singapore, Dec.\ 2023, pp.~12655--12669.
DOI: 10.18653/v1/2023.emnlp-main.778.


\end{thebibliography}

\end{document}
